\documentclass[11pt]{article}
\usepackage[utf8]{inputenc}
\usepackage[T1]{fontenc}
\usepackage{lmodern}
\usepackage[margin=1in]{geometry}
\usepackage{csquotes}
\usepackage[style=authoryear,backend=biber,sorting=nyt,maxbibnames=99,giveninits=true,uniquename=init]{biblatex}
\usepackage[hidelinks]{hyperref}
\addbibresource{references.bib}
\begin{document}
\section*{Controlled Drainage Effects on Flow and Losses}

\subsection{Overview}

A large and internally consistent body of evidence establishes that the way a
subsurface drainage system is designed and operated governs both the volume of
water it discharges and the partitioning of that water among competing
pathways---drain flow, surface runoff, evapotranspiration, and deep
seepage---and that these hydrologic shifts in turn control the losses of
nitrogen, phosphorus, sediment, and salts to receiving waters. Three design
levers recur throughout the literature: drain spacing, drain depth, and outlet
(water-table) control. Controlled drainage (CD), also termed drainage water
management (DWM), raises the outlet elevation to hold a higher water table
during non-cropping or low-demand periods, thereby storing additional water in
the field and reducing cumulative drain outflow \parencite{Singh2020,Strock2011,Blann2009}.
The dominant, well-supported mechanism is a coupling between flow and load:
because most nitrate leaves drained fields dissolved in drain water, reductions
in nitrate-N loss under controlled or shallower drainage are of roughly the same
magnitude as the reductions in outflow, with concentration changes secondary
\parencite{Wesström2007,Frankenberger2024,Lalonde1996,Tolomio2018}.

The headline magnitudes are quantitatively convergent across meta-analyses and
syntheses. Across 15 peer-reviewed studies, conventional drainage volumes
exceeded controlled drainage volumes by an average of 40\%, and controlled
drainage reduced nitrogen loss by 40--50\% and phosphorus loss to surface water
by 25--35\% \parencite{Strock2011}. A meta-analysis of 61 studies (1050 data
pairs) found that controlled drainage decreased drainage volume by 19.2\%,
reduced NO$_3$-N concentration and load, and produced a slight overall crop
yield increase of 0.11\%, framing it as a ``win--win'' practice
\parencite{Wang2020c}. A synthesis of drainage water management drawing on 290
plot-years and 31 sites reported a mean NO$_3$-N load reduction of
13.3~kg~ha$^{-1}$~yr$^{-1}$ (46\% removal efficiency) and a mean flow reduction
of 126~mm~yr$^{-1}$, with dissolved reactive and total phosphorus loads reduced
by 34\% and 36\%, respectively, and with the paired concentration reduction not
significant---confirming that load reduction arises mainly from reduced flow
\parencite{Frankenberger2024}. The MANAGE Drain Load database, compiling 91
publications and 1279 site-years spanning 1961--2012, further anchors the
baseline: the mean and median fractions of precipitation appearing as drainage
were 25\% and 20\%, and wet years produced significantly greater discharge and
nutrient loads than dry years \parencite{Christianson2015}. The evidence base
spans multi-year paired-field experiments, controlled tank and lysimeter
studies, and simulation with DRAINMOD, HYDRUS, SWAP, FLUSH, and SWAT-family
models, giving both process understanding and quantified outcomes across humid,
arid, and irrigated settings.

\subsection{Controlled Drainage and the Reduction of Drain Outflow}

The most robust and frequently replicated finding is that controlled drainage
reduces cumulative and annual subsurface drain discharge relative to free
(conventional) drainage. Field experiments report reductions that, while
site-specific, are consistently large. In Sweden, controlled drainage produced
much lower total drain outflow in both measured years, with the largest
reduction in the second year, and lowered nitrate and phosphorus losses in
proportion to the reduced outflow \parencite{Wesström2001,Wesström2007}. In
Quebec, controlled water-table treatments at 0.25 and 0.50~m above the drain
reduced total drain flow by 58.7\% and 65.3\% in 1992 and by 40.9\% and 95\% in
1993, with nitrate-loss reductions of comparable magnitude (up to 95.7\%), the
net benefit attributed primarily to decreased drain flow
\parencite{Lalonde1996}. Iowa comparisons found controlled and shallow drainage
reduced subsurface drainage volume by 37\% and 46\% (four-year averages of 23
and 20~cm versus 37~cm for conventional), and over a five-year study controlled
and shallow drainage cut annual subsurface flow by 60\% and 58\%
\parencite{Helmers2012,Schott2017}. Ohio field data showed 16--24\% reductions
in annual discharge and a statistically significant drop in mean monthly
discharge from 4.3 to 3.1~cm \parencite{Pease2017}, and across six northwest
Ohio sites the mean reduction in daily managed-zone drainage volume ranged from
40\% to 100\% (with no discharge at one site), the linear mixed models
explaining 57--92\% of the variance \parencite{Gunn2015}. Indiana paired-field
work found the strongest reductions in the dormant season (up to 100\%), with
annual nitrate-load reductions of 15--31\% \parencite{Adeuya2012}, and event
analysis showed controlled drainage reduced event drainage volume and peak flow
and increased time to peak, with antecedent conditions rather than precipitation
characteristics controlling controlled-quadrant response \parencite{Lahdou2019}.

Similar reductions recur worldwide. Italian water-table management reduced total
water discharge by 81\% \parencite{Tolomio2018}; Egyptian and Iranian studies
reported significantly reduced outflow under controlled drainage in semiarid
conditions \parencite{Wahba2001,Jouni2018}; and Indian saline Vertisols showed
controlled drainage cutting drainage water outflow by about 64\% (11.0 versus
30.2~cm) and average drain discharge from 3.10 to 1.12~mm~day$^{-1}$
\parencite{Karegoudar2019}. Michigan simulations projected controlled drainage
reducing annual discharge by 59--72\% depending on management aggressiveness and
period \parencite{BinShokrana2023}, and a Michigan paired-field saturated-buffer
study attributed most performance to the controlled-drainage component, which
reduced annual discharge by 47.1\% and nitrate-N load by 48.0\%
\parencite{Ghane2025}. In cultivated peat soils in Finland, controlled drainage
reduced drain discharge by 862~mm (thin peat) and 457~mm (thick peat), a strong
discharge response accompanied by only a modest groundwater-level rise
\parencite{Tähtikarhu2025}; other Finnish work reported 11--23\% reductions in
mean annual drain discharge and 17--36~cm higher water tables under controlled
drainage \parencite{Salla2022}. Dutch reviews reported reductions of 14\% in a
lysimeter and 16--85\% across clay, silty clay, and sandy loam soils, together
with increased water retention \parencite{deWit2022}. Tank experiments confirmed
the mechanism under strictly controlled boundaries: controlled drainage at 10
and 20~cm reduced cumulative drainage volume by 8\% and 6\% and extended
drainage duration by 17~h \parencite{Zheng2025}. Collectively, these results
establish outflow reduction, extended drainage duration, lower and later peak
flows, and increased in-field storage as consolidated, transferable effects of
outlet control.

\subsection{Nitrate and Phosphorus Loss Reductions and the Flow--Load Coupling}

Because nitrate is highly mobile and travels dissolved in drain water, the
reduction in nitrate-N load under controlled drainage tracks the reduction in
outflow closely. Swedish studies made this explicit: the relative decrease in N
and P loading in drain outflow was of the same magnitude as the reduction in
outflow rate, with high-risk N-loss periods coinciding with high outflow rates
and high soil mineral N \parencite{Wesström2007}. Global syntheses reinforce
that nutrient reduction is due mainly to reduced flow rather than consistent
concentration reductions \parencite{Frankenberger2024,Wang2020c}. Reported
load reductions are large and consistent: 61\% (controlled) and 49\% (shallow)
in Iowa \parencite{Schott2017}; 36\% (controlled) and 29\% (shallow) in Iowa
\parencite{Helmers2012}; about 47\% at the watershed scale in Indiana, where
drain flow fell from 11.0 to 5.9~cm and nitrate load from 236.3 to 126.4~ton,
with most reduction occurring in the November--March winter control period
\parencite{Ale2012}; and 54\% for the saturated-buffer system, dominated by the
controlled-drainage component \parencite{Ghane2025}. In India, controlled
drainage roughly halved seasonal NO$_3$-N loss (from 11.2 to
5.54~kg~ha$^{-1}$ per season) \parencite{Karegoudar2019}, and Italian water-table
management lowered annual N loss by 92\% (from 29 to 2~kg NO$_3$-N~ha$^{-1}$) and
P loss by 65\%, with cumulative nitrate loss strongly correlated with cumulative
outflow volume \parencite{Tolomio2018}. Chinese field and modeling studies found
controlled drainage increased soil and crop NO$_3$-N while reducing leaching and
drain nitrate loss by roughly 15--29\% \parencite{Dou2021,Dou2023}, and Iranian
work found lower flow-weighted nitrate concentrations and nitrate losses under
controlled drainage \parencite{Jouni2018}. A broad review of edge-of-field and
in-field practices placed controlled drainage among effective mitigation
options, with reported NO$_3$ loss reductions of 18--75\% for controlled
drainage, alongside 12--100\% for bioreactors and 25--78\% for wetlands
\parencite{Castellano2019}.

Denitrification provides a secondary but well-documented loss-reduction
mechanism where the water table is held high. Canadian farm-scale work found a
greater relative abundance of nitrate-reducing/facultative bacteria under
controlled drainage, supporting denitrification as a nitrate-loss pathway, while
also reducing total suspended solids and particulate phosphorus by damping
macropore flow \parencite{Tanga2025}. However, several syntheses caution that in
poorly drained soils the nitrate reduction is attributable mainly to reduced
tile effluent rather than to increased denitrification \parencite{THOMAS1995},
and Albanian field data found nutrient concentrations were not strongly affected
by drainage-system type, with outflow volume the dominant control on total
nutrient transport \parencite{Grazhdani1996}. Modeling in Belgium reproduced an
adverse relationship between denitrification and leaching---wider spacings and
higher water tables increase denitrification and reduce leaching but also reduce
plant uptake---so that total winter NO$_3$-N losses could be substantially
reduced under controlled drainage \parencite{ElSadek2002,ElSadek2002b}.

Phosphorus responses are generally favorable but smaller and occasionally
reversed. Controlled drainage lowered P loads in Sweden and Iran
\parencite{Wesström2007,Wesström2001,Jouni2018}, and the DWM synthesis reported
DRP and TP load reductions of about 34\% and 36\%
\parencite{Frankenberger2024}. In contrast, North Carolina simulations found
controlled drainage produced a small increase in phosphorus efflux even as it
cut nitrate efflux by as much as 34\% \parencite{DEAL1986}, and Iranian data
showed slightly higher P concentrations under controlled drainage though lower P
losses overall \parencite{Jouni2018}. Broader syntheses note that timing matters:
the most hazardous periods for surface-water pollution are autumn--winter, when
rainfall coincides with fertilizer on bare soil \parencite{Tolomio2018}, and P
peaks can shift from late winter under conventional drainage to periods of
highest outflow under controlled drainage \parencite{Wesström2007}. A recurring
caveat is that controlled drainage may increase nutrient loss through seepage
and surface runoff, a trade-off that partially offsets drain-flow reductions
\parencite{Frankenberger2024}.

\subsection{Repartitioning of Water Among Pathways}

A consolidated finding is that controlled and shallow drainage do not eliminate
the stored water but redirect it, most consistently into increased surface
runoff, and secondarily into evapotranspiration and seepage. DRAINMOD
simulations for Iowa showed that at a drainage intensity of 4.0~mm~day$^{-1}$,
shallow and controlled drainage each reduced average annual subsurface drainage
by about 25--27~mm (roughly 15\%), with most of that reduction reappearing as
increased surface runoff and little change in evapotranspiration
\parencite{Singh2007b}. Belgian, North Carolina, and Canadian studies
consistently report that controlled drainage decreases subsurface drainage and
increases surface runoff \parencite{ElSadek2002,ElSadek2002b,DEAL1986,Skaggs1995,Golmohammadi2016,Ndulue2024}.
At the watershed scale in Ontario, controlled drainage decreased subsurface
outflow and increased surface runoff, most often during the March--April
snowmelt period, while producing a shallower water table and only a slight
increase in evapotranspiration \parencite{Golmohammadi2016}. Polish simulations
projected controlled drainage and reduced subsurface outflow accompanied by an
approximately fivefold increase in surface runoff from drained arable land
\parencite{Sojka2020}. Finnish modeling similarly found controlled subsurface
drainage alone reduced drain discharge the most, whereas open-ditch damming
alone increased drain discharge and reduced groundwater outflow, and all
controlled-drainage methods increased topsoil-layer runoff while raising
evapotranspiration and reducing total runoff overall
\parencite{Isomäki2024}. Drainage-design choices are most clearly expressed in
the runoff outflow pathways \parencite{Isomäki2024}. Increased evapotranspiration
under controlled drainage---up to a 10\% ETact increase for corn---is a
recognized secondary sink that reduces total water lost to drainage
\parencite{deWit2022,DEAL1986,Ndulue2024}.

At the field-water-balance scale, tile drainage dominates the outflow partition
in humid Corn Belt settings. In Indiana, surface runoff represented only 1--10\%
of annual precipitation while subsurface drainage represented 26--45\%, with a
simulated balance of 7\% runoff, 34\% drain flow, and 51\% evapotranspiration,
and all ponding events generated by saturation excess rather than infiltration
excess \parencite{Saadat2020}. Iowa pothole modeling similarly found tile
discharge and evapotranspiration accounting for 52.4\% and 47.5\% of
precipitation, with drainage reducing but not eliminating ponding
\parencite{Amado2018,Martin2019}. In UK field-drained catchments, subsurface
field drains were the dominant pathway transferring runoff, sediment, and---the
authors emphasize---most sediment to the stream, with drainage responding in
every monitored event whereas surface runoff occurred in only some
\parencite{Deasy2009}. These partitions underscore that in artificially drained
landscapes the drain, not the soil surface, is usually the principal conduit.

Relative to undrained or predominantly surface-drained conditions, however,
improved subsurface drainage generally reduces surface runoff, peak outflow
rates, sediment, and phosphorus even as it increases nitrate delivery
\parencite{FAUSEY1995,SKAGGS1994,THOMAS1995,Gramlich2018}. In Oregon, tile
drainage lowered perched water tables and reduced watershed runoff and sediment
yield by about 65\% and 55\%, eliminated baseflow, increased lag time, and
lowered peak and average flow \parencite{ISTOK1983}. Albanian sites recorded
annual surface-runoff reductions of 30--56\% and reduced soil loss under
subsurface drainage, though total nitrogen losses rose 22--25\%
\parencite{Grazhdani1996}. Minnesota watershed modeling found tile drainage
reduced surface runoff, soil loss, and sediment yield, especially under no-till
\parencite{Maalim2013}. In North Carolina, improved subsurface drainage reduced
surface runoff by 192~mm~yr$^{-1}$ (66\%) on a Wasda soil while raising total
outflow by 10\%, and attenuated peak runoff strongly at the field scale (from
101 to 28~mm~day$^{-1}$ for a three-year storm) with differences persisting to
the 6216~ha watershed outlet \parencite{KONYHA1992}. Sandy loam soils in Canada
showed subsurface drainage minimizing surface runoff and total phosphorus loss,
whereas the reverse held on clay loam, and phosphorus speciation was mainly a
function of soil texture rather than transport pathway
\parencite{Eastman2010}. Mechanistically, good subsurface drainage increases
infiltration and drainage rates during rainfall and lowers solute concentrations
in surface runoff \parencite{Tong2016,Tao2017}, and by lowering the water table
it limits saturation-excess runoff and surface sealing \parencite{Augeard2005}.

\subsection{Drain Depth and Spacing as Design Levers}

Beyond outlet control, spacing and depth systematically shift the flow--loss
balance, and the evidence is coherent across field and simulation studies.
Increasing drain spacing reduces subsurface drainage and nitrate leaching to
drains but increases surface runoff, denitrification, and nitrate losses in
runoff, while decreasing plant uptake---defining an optimum-spacing trade-off
between crop production and environmental impact \parencite{ElSadek2002,ElSadek2002b}.
DRAINMOD-N simulations in North Carolina found that increasing spacing from 20
to 40~m decreased predicted NO$_3$-N losses by 47\% while still meeting
production goals \parencite{Skaggs1995}, and the MANAGE database confirmed that
wider spacing and shallower placement tend to reduce dissolved N loss, with
deeper drains showing disproportionately higher N loads at a given discharge
\parencite{Christianson2015}. Shallower drains, correspondingly, move less water
and nutrient offsite: reducing drain depth from 120 to 90~cm produced roughly
20\% reductions in annual drainage volume and nitrate-N over six years in
Minnesota (and in a 90-year simulation), and a 76~versus~120~cm comparison
yielded about 40\% reductions \parencite{Strock2011}. Regional syntheses agree
that shallow drains contain less sediment and nutrients than deeper drains and
that narrow spacing moves more water, nutrients, and pesticides offsite than
wide spacing \parencite{FAUSEY1995}. Chinese salinity studies identified
site-specific optima---8~m spacing at 0.7~m depth for a greenhouse soil
\parencite{Shao2012} and 0.9~m depth with 10~m spacing for chloride removal
\parencite{Yan2013b}.

Because different spacing--depth combinations can achieve similar drainage
intensity, design can be tuned to jointly meet productivity and water-quality
goals. Latvian field experience over 16--17 years on about 200 fields showed
that similar drainage intensity can be provided by different combinations of
spacing and depth, favoring the most economical choice, and that deep systematic
drainage (1.3--1.5~m) gave the best draining and economic indices in sandy and
loamy soils \parencite{SHKINKIS1979}. Belgian economic optimization identified a
conventional 50~m spacing at 1.25~m depth as the profit optimum, but showed that
widening spacing from 50 to 100~m reduced total nitrate-N loss by 20\% (from 37
to 29~kg~ha$^{-1}$) at only a small profit cost, and that different
depth--spacing scenarios can yield similar nitrate-loss reductions
\parencite{ElSadek2002b}. North Carolina simulations found deeper drains
increased profit but substantially increased NO$_3$-N losses, while improved
surface drainage allowed wider spacings and lower nitrate losses; the general
prescription is to use the minimum subsurface intensity required for production,
fit systems to soil properties, and apply controlled drainage when maximum
drainage is unnecessary \parencite{Breve1998}. Irrigated-system reviews likewise
conclude that deep drains are often unnecessary for salinity control and that
shallower drains with better operational management can reduce design discharges
without sacrificing yield \parencite{Ritzema2011}.

\subsection{Shallow, Multi-level, and Enhanced Drainage Configurations}

Shallow drainage emerges as a design that reproduces much of the flow-- and
load-reduction benefit of controlled drainage. Iowa studies found shallow
drainage reduced subsurface volume by 46--58\% and nitrate loads by 29--49\%,
generally without reducing crop yield (unlike controlled drainage, which
sometimes did) \parencite{Helmers2012,Schott2017}. A multi-state eastern-USA
analysis showed shallow drains produced quicker water-table drawdown, reduced
year-to-year corn-yield variability, and cut total flow by 3.4~cm versus deep
drains (with discharge reduction intensifying with annual precipitation), though
they increased surface runoff and had a lower benefit--cost ratio and longer
payback than deep drains \parencite{Ghane2021b}.

In saline irrigated settings, drain depth strongly controls the salinity of
exported water. Australian experiments found deep drains flowed continuously and
exported the largest salt loads (5867~kg~ha$^{-1}$), whereas shallow drains
flowed only after irrigation or rainfall, had low salinity, and produced the
greatest salt-load reduction while still providing adequate waterlogging control
\parencite{Christen2001b}. Multi-level (bilevel) systems combining deep and
shallow drains improved on single-level systems: they cut the time the water
table exceeded 0.5, 0.75, and 1.0~m by 84\%, 67\%, and 41\%, leached the root
zone faster, and produced much lower shallow-drain salinity (4.7 versus
31.7~dS~m$^{-1}$) while removing similar total salt
\parencite{Hornbuckle2007}. Analytical theory shows that adding shallow drains
between deep drains raises system drainage rate, produces shallower and more
evenly distributed flow paths, and allows wider deep-drain spacing for equivalent
total drainage \parencite{Hornbuckle2012}. Managed deep drains occupy an
intermediate position, reducing volume and salinity relative to unmanaged deep
drains \parencite{Ayars2006b,Christen2001b}.

Engineered enhancements to the drain itself further alter flow. ``Improved''
subsurface drainage that adds a permeable filter increased discharge---about 1.9
times conventional discharge at equal ponding depth---accelerated water-table
drawdown (over 15\% faster to reach 40~cm), and reduced surface runoff more
effectively, especially under high-intensity rainfall
\parencite{Tao2016,Tao2017}. Combined open-ditch plus subsurface designs and
improved filters increased flow rates and flooding/waterlogging removal
efficiency, with medium--medium and shallow--deep combinations offering good
performance at lower construction difficulty \parencite{Ren2021,Tao2025}.
Siphon-type subsurface drainage raised flow rates by roughly 27--90\% over
conventional and improved variants and shortened flooding and waterlogging
removal times substantially, verified against HYDRUS-2D
\parencite{Ren2023}. Artificial macropore flow tubes generated drain flow where
none occurred in a base treatment, increased drainage volume with tube density,
and promoted a more uniform, desalinated soil profile \parencite{Chen2021}.
Supplementary drainage of a clay field increased mean annual drain discharge by
a factor of 2.6, decreased topsoil-layer runoff by 75\% and groundwater outflow
by 18\%, and improved spring drainage efficiency without affecting
evapotranspiration \parencite{Häggblom2019}. Controlled drainage combined with
subirrigation raised groundwater levels, supplied 170--410~mm of infiltration,
and increased yields by up to 64\% in Dutch pilots, though supplies were
insufficient in extreme dry seasons \parencite{deWit2022,deWit2024}, and
subirrigation reduced nitrate-N loading and concentration in earlier North
American work \parencite{FAUSEY1995,Almen2021,Madramootoo2007}. Land retirement
with drains closed can eliminate drainage entirely, with downgradient retirement
generating larger reductions than upgradient \parencite{Purkey2001}.

\subsection{Salinity, Salt Load, and Trade-offs in Irrigated Systems}

In arid and irrigated regions the flow--loss coupling extends to salts, and
controlled drainage again reduces exported volumes and loads but introduces an
in-field salinity trade-off. Controlled drainage reduced total drainage outflow
and salt load in the United States and Australia \parencite{Ayars2006b}, cut
cumulative salt discharge by 5--7\% in tank experiments (while raising drainage
water salinity, especially late) \parencite{Zheng2025}, and reduced salt loads
in Iranian managed paddy systems, where free drainage produced salt loads 4--6
times higher than managed treatments \parencite{DarziNaftchali2018b}. Indian
saline-Vertisol trials found controlled drainage saved about 17\% of irrigation
water and roughly halved nitrogen loss, but reclaimed soil salinity more slowly
than conventional drainage \parencite{Karegoudar2019}, and Iranian work reported
soil salinity decreasing under free drainage but increasing slightly under
controlled drainage---an explicit productivity--salinity trade-off
\parencite{Jouni2018}. Reviews of Australian systems highlighted that many
subsurface systems removed more water than needed, mobilizing stored subsoil
salt, and that managing horizontal drains to avoid full drawdown could reduce
salt loads significantly \parencite{Christen2001}. Egyptian controlled drainage
reduced outflow, nitrate, and phosphorus losses in semiarid conditions, with
potential to lower the cost of pumping drainage water to the sea
\parencite{Wahba2001}.

The counterpart consolidated finding is that subsurface drainage is a proven,
economically favorable technology for reclaiming waterlogged and saline
irrigated land. Indian pilot areas recorded rice-yield increases of 30--120\%
and average increases of 54--136\% across crops, with cost--benefit ratios of
2.5--4.6 and payback periods of 2--3 years \parencite{Satyanarayana2007}, and
long-term monitoring found water tables 25\% deeper and soil salinity 50\% lower
in drained fields, with root-zone salinity falling from 11.5 to 6.0~dS~m$^{-1}$
and favorable cost--benefit ratios of 1.2--3.2 \parencite{Ritzema2008b}. Deep
drains proved unnecessary for salinity control, with shallow drains (0.8--1.0~m)
sufficiently effective \parencite{Ritzema2011,Satyanarayana2007}. Egyptian
drainage lowered water tables to 0.8--1.0~m within days of irrigation and
reduced saline soils from 40\% to 5--8\% (new non-saline areas) within four
years, proving highly profitable \parencite{Ali2001}. Comparable reclamation
successes and adequate water-table and salinity control were documented in
Mexico, Turkey, and China \parencite{MartínezBeltrán2007,Güngör2023,Shao2012,Yan2013b}.
Italian analysis showed drainage altered redox conditions above the drains,
favoring nitrification and confining active water circulation to the zone above
the drains \parencite{Mastrocicco2013}. At the basin scale, however, persistent
waterlogging and salinity in Central and South Asia reflect deteriorated
infrastructure and inadequate saline-effluent disposal, underscoring that
drainage must be managed within an integrated salt balance
\parencite{Dukhovny2007,Qureshi2008,Ritzema2016,SKAGGS1994}.

\subsection{Crop Yield, Water Use, and Agronomic Outcomes}

The agronomic evidence supports the conclusion that controlled and shallow
drainage can reduce losses while maintaining---or modestly improving---yield and
water productivity. The 61-study meta-analysis found an overall yield increase of
0.11\%, with dryland crops increasing slightly and paddy rice decreasing
\parencite{Wang2020c}, and North Carolina studies indicated an average yield
increase of about 5\% when controlled drainage is well managed
\parencite{Strock2011}. Dutch reviews reported yield gains of 2--5\% for corn,
with one instance of a 3\% decrease when the water table was too shallow early
in the season \parencite{deWit2022}. Controlled drainage improved nitrogen
efficiency in Sweden through lower drain N loads and higher crop uptake
\parencite{Wesström2007}, increased yield and water-use efficiency in Iranian
and Chinese systems \parencite{Jouni2018,Dou2021,Dou2023}, and increased water
productivity and yield in Iranian managed paddy fields with quick payback
\parencite{DarziNaftchali2018b}. Not all yield effects are positive: an Iowa
comparison found controlled drainage produced a significant yield reduction
relative to conventional drainage even as shallow drainage did not
\parencite{Helmers2012}, whereas a five-year Iowa study found yields similar
among conventional, controlled, and shallow drainage, with only the undrained
treatment reducing corn and soybean yields by 6\% \parencite{Schott2017}.

Drainage per se---versus no drainage---consistently stabilizes and improves
agronomic performance. Iowa nitrogen-response work found controlled drainage had
the highest agronomic efficiency at the agronomic optimum N rate and the least
interannual variability, while no drainage had the highest optimum N rate, the
greatest N surplus, and the most variability \parencite{Maas2022}. A conceptual
analysis showed drainage in continuous maize lowered the optimum N rate from 214
to 189~kg~N~ha$^{-1}$, raised zero-N yield by 10\% and agronomic efficiency by
14\%, and reduced global warming potential by 56\% through lower fertilizer and
N$_2$O emissions, positioning controlled and shallow drainage as designs that
reduce leaching while maintaining yield and conserving soil water and carbon
\parencite{Castellano2019}. Drainage also advances field trafficability: Canadian
work found subsurface drainage advanced spring trafficability by 2--5 weeks and
recommended a 30.5~m spacing, with intensity mattering most at wide spacing
\parencite{PAUL1983}, and Iowa analysis showed drained treatments enabled earlier
target planting dates than undrained soils in wet years, with controlled and
shallow drainage reducing nitrate loss without delaying planting or reducing
yield \parencite{Schott2017b}. Managed subsurface drainage accelerated ponded-
and groundwater removal and enabled timely second-crop establishment in
Bangladesh, India, and Iran \parencite{Okwany2016,Prathapar2018,Rahimi2024}.

\subsection{Persistence of Benefits and Behavior under Climate Change}

Climate-projection studies consistently conclude that the hydrologic and
water-quality benefits of controlled drainage persist, and in some seasons
strengthen, under future conditions. In Ohio, conventional discharge was
projected to decline in all seasons (most in autumn), while controlled-drainage
annual discharge reduction remained near 4--6\% and summer growing-season
reductions increased, indicating greater warm-season value
\parencite{Pease2017}. Michigan simulations projected free-drainage discharge
rising 20\% by mid-century while controlled drainage continued to reduce
discharge by 63--72\%, with future and historical performance statistically
indistinguishable \parencite{BinShokrana2023}. Finnish future scenarios projected
controlled drainage effects on groundwater and discharge remaining similar to the
historical interval, with annual drain-discharge reductions of roughly
21--46~mm depending on management \parencite{Salla2022}. Canadian projections of
warmer, wetter conditions found controlled drainage reducing subsurface drainage
more strongly than free drainage and raising evapotranspiration, though rising
temperature drove dry-stress yield reductions, motivating capture, storage, and
reuse of drainage water as an adaptation \parencite{Ndulue2024}. Chinese and
Polish studies found that appropriate surface-depression and outlet-blocking
strategies (e.g., starting regulation on 1~March) could reduce future runoff and
maintain groundwater at present-day levels \parencite{Awad2021,Sojka2020}.
Field monitoring in Norway demonstrated that subsurface systems react rapidly to
precipitation-induced groundwater rise and can handle high precipitation amounts
similar to climate-change projections, with two-exponential recession behavior
capturing fast and slow flow \parencite{Deelstra2015}, and Polish work quantified
seasonal retention and outflow factors, showing the greatest flood risk under
rapid snowmelt on near-saturated soil \parencite{Szejba2020}.

\subsection{Landscape Extent, Monitoring, and the Evidence Base}

Quantifying where and how much land is drained frames the significance of these
flow and loss effects. Regional mapping in Germany estimated drained proportions
of about 25\% overall, highest on hydromorphic groundwater- and
stagnant-water-influenced soils (up to 30--40\%) and lowest on leachate soils,
with the method transferable to river-basin and national scales
\parencite{Hirt2005,Hirt2011}. Multiple United States drainage-extent datasets
diverge substantially at the watershed scale, indicating that the best estimate
depends on local drainage conditions and dominant practice
\parencite{Jame2022}. Regional drainage-requirement and monitoring studies in
India, Turkey, Egypt, and elsewhere documented drainable-surplus estimates,
recommended design rates (e.g., 2~mm~day$^{-1}$), and generally good system
performance, with water tables dropping below the root zone within days of
irrigation \parencite{Kumar2006,Güngör2023,Ali2001,MartínezBeltrán2007,Rimidis2003}.
Setback analyses show drain depth is the key determinant of lateral setback
distance for protecting wetlands \parencite{Tangen2018}, and vulnerability
indices require distinguishing surface from subsurface artificial drainage to
avoid mis-classifying runoff and leaching risk \parencite{Baffaut2020}. Runoff
from small catchments and drainage fields agreed well, supporting regional
prediction \parencite{Veinbergs2022}, while watershed models such as SWAT were
found to under-represent artificial subsurface drainage processes without
explicit adaptation \parencite{Manguerra1998,Golmohammadi2016}. Broadly,
subsurface drainage relative to surface drainage lowers sediment, phosphorus, and
pesticide concentrations but carries much greater nitrate loads, and its
hydrologic and water-quality effects vary widely with soil, topography, land use,
and system design---so conservation designs such as controlled drainage offer
substantial potential to reduce nitrate losses and improve water quality where
appropriately sited \parencite{Blann2009,Gramlich2018,FAUSEY1995,SKAGGS1994}.
Socio-economic and landscape considerations temper further expansion in some
contexts: Finnish analyses found replacing remaining open ditches with subsurface
drainage could reduce social welfare through biodiversity and landscape losses,
so surface drainage remained socially optimal for lower land qualities
\parencite{HietalaKoivu2004,Saikkonen2014}. The overarching, well-supported
conclusion across this literature is that drainage systems can be deliberately
designed and managed---through spacing, depth, and outlet control---to satisfy
agricultural production and environmental objectives simultaneously
\parencite{Skaggs1995,Breve1998,ElSadek2002b,Strock2011,Ritzema2016}.

\printbibliography
\end{document}
